\chapter{تحويل فورييه}\label{ch:b3:fouriertransform}

تفكّك متسلسلات فورييه الإشاراتِ الدورية إلى توافقيات متقطّعة؛
ويفعل \emph{تحويل} فورييه الشيء نفسه بالإشارات على المستقيم
كله، بمتصلٍ من الترددات. فهو يحوّل الاشتقاق إلى ضرب، والالتفاف
إلى جداء، والدوال الغاوسية إلى دوال غاوسية --- ولهذا يحل
المعادلات التفاضلية، ويقود معالجة الإشارة، وسيبرهن على مبرهنة
النهاية المركزية في \cref{ch:b3:clt}. ويطوّر هذا الفصل نظرية $L^1$
(ريمان--لوبيغ، والقلب، والتباين)، \omterm{def:b3:fouriertransform:schwartz}{وصنف شوارتز} حيث يكون التحويل
تقابلًا تامًّا، ونظرية $L^2$ (بلانشرِل: أي إن التحويل، إلى غاية
ثابت، مؤثر واحدي)، مع تطبيقين بارزين: معادلة الحرارة، محلولةً
من أولها إلى آخرها في مسألة نهاية الأسبوع، وصيغة بواسون في
الجمع. والاصطلاح:
\[
\hat f(\xi) = \int_\R f(x)\,\eu^{-\iu\xi x}\,\dd x .
\]

\section{التحويل على \texorpdfstring{$L^1$}{L1}}

\begin{proposition}\label{prop:b3:fouriertransform:basic}
من أجل $f \in L^1(\R)$: تكون $\hat f$ معرَّفة جيدًا، ومحدودة
($\norm{\hat f}_\infty \leq \norm f_1$)، ومتصلة، و:
\begin{enumerate}
  \item $\widehat{\tau_af}(\xi) = \eu^{-\iu a\xi}\hat f(\xi)$ و $\widehat{\eu^{\iu ax}f}(\xi) = \hat f(\xi - a)$؛
  \item $\widehat{f(\cdot/\lambda)}(\xi) = \lambda\hat
        f(\lambda\xi)$ من أجل $\lambda > 0$؛
  \item إذا كان $xf \in L^1$، فإن $\hat f$ من الصنف
        $\mathcal C^1$ مع $(\hat f)'(\xi) = \widehat{(-\iu x)f}(\xi)$؛
  \item إذا كان $f \in \mathcal C^1$ مع $f' \in L^1$ (ومع $f \to
        0$ عند
        $\pm\infty$، وهو تلقائي هنا)، فإن $\widehat{f'}(\xi) = \iu\xi\hat f(\xi)$؛
  \item $\widehat{f * g} = \hat f\,\hat g$ من أجل $f, g \in L^1$.
\end{enumerate}
\end{proposition}

\begin{proof}
أما الحد: $\abs{\hat f} \leq \int\abs f$. وأما الاتصال: فبالتقارب المهيمن مع المهيمِن
$\abs f$ (\cref{thm:b3:lebesgue:paramcont}). و(1) و(2): بالتعويضات (\cref{thm:b3:product:linearchange}).
و(3): بالاشتقاق تحت علامة التكامل، بالمهيمِن $\abs{xf}$
(\cref{thm:b3:lebesgue:paramdiff}). و(4): أولًا، للمقدار $f(x) = f(0)
+ \int_0^xf'$ نهايةٌ عند $\pm\infty$
(لأن $f' \in L^1$)، ولا بد أن تكون $0$ (لأن $f \in L^1$)؛ ثم
نكامل بالتجزئة على $[-A,
A]$ ونجعل $A \to \infty$. و(5): بفوبيني،
وهو مشروع لأن $(x,y)\mapsto f(x - y)g(y)\eu^{-\iu\xi x}$ قابل للمكاملة مطلقًا (\cref{thm:b3:product:convolution}):
\[
\widehat{f*g}(\xi) = \iint f(x - y)g(y)\eu^{-\iu\xi(x - y)}
\eu^{-\iu\xi y}\dd x\,\dd y = \hat f(\xi)\,\hat g(\xi).
\]
\end{proof}

\begin{example}\label{ex:b3:fouriertransform:gaussian}
الدالة الغاوسية: من أجل $a > 0$،\index{الدالة الغاوسية، تحويل
فورييه}
\[
\widehat{\eu^{-ax^2}}(\xi) =
\sqrt{\frac\pi a}\;\eu^{-\xi^2/4a} :
\]
حسب \cref{exo:b3:lebesgue:7} (بحيلة المعادلة التفاضلية $F' = -\frac\xi{2}F$، معادةَ
التحجيم)، أو حسب (3): إذ تحقق $g = \widehat{\eu^{-ax^2}}$ أن $g'(\xi) = -\frac{\xi}{2a}g(\xi)$ (بالمكاملة
بالتجزئة)، و $g(0) = \sqrt{\pi/a}$. فالدوال الغاوسية نقاطٌ صامدة للتحويل إلى
غاية تحجيم --- وهذا هو السبب العميق في سيادتها على مبرهنة
النهاية المركزية.
\end{example}

\begin{theorem}[ريمان--لوبيغ]
\label{thm:b3:fouriertransform:riemannlebesgue}
من أجل $f \in L^1(\R)$: $\hat f(\xi) \to 0$ حين $\abs\xi \to
\infty$. ومن ثَمّ $\widehat{\phantom f} \colon L^1 \to \mathcal
C_0(\R)$ (أي
الدوال المتصلة المنعدمة عند اللانهاية).\index{المبرهنة المساعدة
لريمان--لوبيغ}
\end{theorem}

\begin{proof}
من أجل دالة مميّزة لفترة، $\widehat{\mathbf 1_{\intcc
ab}}(\xi) = \frac{\eu^{-\iu a\xi} - \eu^{-\iu b\xi}}{\iu\xi}
\to 0$؛ ومنه من أجل الدوال السلّمية.
والدوال السلّمية كثيفة في $L^1$ (\cref{thm:b3:lp:density}(1) مع تقريب
المجموعات المنتهية \omterm{def:b3:measure:measure}{القياس} باتحادات منتهية لفترات، \cref{exo:b3:measure:7})،
والتحويل متصل بالمعنى $\norm\cdot_\infty$--$\norm\cdot_1$: فمن أجل $\norm{f -
s}_1 < \varepsilon$،
$\limsup_{\abs\xi\to\infty}\abs{\hat
f(\xi)} \leq \varepsilon$.
\end{proof}

\section{القلب والتباين}

\begin{lemma}[صيغة الضرب]
\label{lem:b3:fouriertransform:multiplication}
من أجل $f, g \in L^1(\R)$: $\displaystyle\int \hat f\,g =
\int f\,\hat g$.
\end{lemma}

\begin{proof}
يساوي الطرفان $\iint f(x)g(\xi)\eu^{-\iu x\xi}\dd
x\,\dd\xi$ (بتونيلي--فوبيني: إذ التكامل المزدوج
للقيمة المطلقة هو $\norm f_1\norm g_1$).
\end{proof}

\begin{theorem}[القلب]\label{thm:b3:fouriertransform:inversion}
لتكن $f \in L^1(\R)$.\index{مبرهنة قلب فورييه}
\begin{enumerate}
  \item (القابلية للجمع الغاوسي) من أجل كل $x$،
        \[
        (f * g_\varepsilon)(x) =
        \frac1{2\pi}\int_\R \hat f(\xi)\,
        \eu^{-\varepsilon\xi^2}\,\eu^{\iu x\xi}\,\dd\xi,
        \qquad\text{حيث} g_\varepsilon(y) =
        \frac{1}{2\sqrt{\pi\varepsilon}}\,
        \eu^{-y^2/4\varepsilon},
        \]
        و $f * g_\varepsilon \to f$ في $L^1$ حين $\varepsilon
        \to 0$.
  \item وإذا كان فوق ذلك $\hat f \in L^1$، فإنه من أجل $x$ في
        كل مكان تقريبًا
        \[
        f(x) = \frac{1}{2\pi}\int_\R \hat
        f(\xi)\,\eu^{\iu x\xi}\,\dd\xi ,
        \]
        وتقبل $f$ ممثّلًا \omterm{def:b3:topology:continuity}{متصلًا}.
  \item (التباين) إذا كان $\hat f = 0$ فإن $f = 0$ في كل مكان
        تقريبًا.
\end{enumerate}
\end{theorem}

\begin{proof}
(1) نثبّت $x$ ونطبّق \cref{lem:b3:fouriertransform:multiplication} على $f$ وعلى $g(\xi) = \frac1{2\pi}\eu^{-\varepsilon\xi^2}\eu^{\iu x\xi}$: فحسب
\cref{ex:b3:fouriertransform:gaussian} (مع قاعدة التضمين)،
\[
\hat g(y) = \frac1{2\pi}\sqrt{\frac\pi\varepsilon}\,
\eu^{-(y - x)^2/4\varepsilon} = g_\varepsilon(x - y),
\]
ومنه $\frac1{2\pi}\int\hat f(\xi)\eu^{-\varepsilon\xi^2}
\eu^{\iu x\xi}\dd\xi = \int f(y)g_\varepsilon(x - y)\dd y =
(f*g_\varepsilon)(x)$. والدوال $g_\varepsilon$ وحدةٌ تقريبية: $g_\varepsilon \geq 0$،
$\int g_\varepsilon = 1$ (بالتكامل الغاوسي)، وهي تتركّز عند $0$؛ وينطبق برهان
\cref{thm:b3:lp:regularization}(2) حرفيًّا (إذ لم يُستعمل سوى $\int g_\varepsilon = 1$ والتركّز: فمن
أجل الذيل، $\int_{\abs y > \delta}g_\varepsilon \to
0$): ومنه $\norm{f * g_\varepsilon - f}_1 \to 0$.

(2) إذا كان $\hat f \in L^1$: فإن الطرف الأيمن في (1) يتقارب،
بالتقارب المهيمن (بالمهيمِن $\abs{\hat f}$)، إلى $\frac1{2\pi}\int\hat
f(\xi)\eu^{\iu x\xi}\dd\xi$ من أجل
\emph{كل} $x$، ودالة النهاية هذه متصلة (بالتقارب المهيمن مرة
أخرى). ومن جهة أخرى $f
* g_\varepsilon \to f$ في $L^1$، ومنه على امتداد متتالية
جزئية في كل مكان تقريبًا (\cref{thm:b3:lp:complete}): فتتوافق النهايتان في كل
مكان تقريبًا.

(3) يجعل $\hat f = 0$ الطرفَ الأيمن في (1) منعدمًا: أي $f *
g_\varepsilon = 0$
من أجل كل $\varepsilon$، و $f *
g_\varepsilon \to f$ في $L^1$: ومنه $f = 0$ في كل
مكان تقريبًا.
\end{proof}

\section{صنف شوارتز}

\begin{definition}\label{def:b3:fouriertransform:schwartz}
\emph{صنف شوارتز}\index{صنف شوارتز} $\mathcal
S(\R)$ مؤلَّف من الدوال $f$
من الصنف $\mathcal C^\infty$ التي تحقق $\sup_x\abs{x^m f^{(n)}(x)} < \infty$ من أجل كل $m, n \geq 0$ (أي إن
جميع مشتقاتها تتناقص أسرع من أي قوة). أمثلة: $\eu^{-ax^2}$،
و $\mathcal C_c^\infty$. ومن الواضح أن $\mathcal S
\subseteq L^p$ من أجل كل $p$ (بالحد بالمقدار
$C(1 + x^2)^{-1}$)، وأن $\mathcal S$ مستقر بالاشتقاق، وبالضرب في كثيرات
الحدود، وبالجداءات.
\end{definition}

\begin{theorem}\label{thm:b3:fouriertransform:schwartz}
يرسل تحويل فورييه $\mathcal S(\R)$ تقابليًّا على نفسه، ومقلوبه
$\check g(x) = \frac1{2\pi}\int
g(\xi)\eu^{\iu x\xi}\dd\xi$.
\end{theorem}

\begin{proof}
لتكن $f \in \mathcal S$. وبتكرار \cref{prop:b3:fouriertransform:basic}(3)، $(\hat f)^{(n)} =
\widehat{(-\iu x)^nf}$ (إذ كل
$x^kf \in L^1$)؛ وبتكرار (4) مع $h = (-\iu x)^nf \in \mathcal S$ (وجميع مشتقاتها \omterm{def:b3:lebesgue:l1}{قابلة
للمكاملة})، $(\iu\xi)^m\hat h =
\widehat{h^{(m)}}$. وبالضمّ،
\[
\abs{\xi^m\,(\hat f)^{(n)}(\xi)}
= \bigl|\widehat{\,h^{(m)}}(\xi)\bigr|
\leq \bigl\|\bigl((-\iu x)^nf\bigr)^{(m)}\bigr\|_1 < \infty
\]
بانتظام في $\xi$: ومنه $\hat f \in \mathcal S$. وبما أن $\hat f \in L^1$، فإن
القلب (\cref{thm:b3:fouriertransform:inversion}(2)) يصح في كل مكان (لأن الطرفين متصلان):
$\check{\hat f} = f$، وبالتناظر $\widehat{\check g} = g$ (إذ التحويل المرافق هو $g \mapsto \frac1{2\pi}\hat g(-\cdot)$، وهو
أيضًا يحفظ $\mathcal S$): فيكون تقابلًا.
\end{proof}

\section{بلانشرِل و \texorpdfstring{$L^2$}{L2}}

\begin{theorem}[بلانشرِل]\label{thm:b3:fouriertransform:plancherel}
من أجل $f \in L^1 \cap L^2(\R)$:\index{مبرهنة بلانشرِل}
\[
\norm{\hat f}_2^2 = 2\pi\,\norm f_2^2 .
\]
ومن ثَمّ يمتدّ $\widehat{\phantom f}$ بكيفية وحيدة إلى تطبيق خطي متصل $\mathcal F \colon L^2(\R) \to L^2(\R)$
يحقق $\norm{\mathcal Ff}_2 = \sqrt{2\pi}\norm f_2$؛ ويكون $\mathcal
F$ تقابلًا، مع $\mathcal F^{-1} =
\frac1{2\pi}\,\mathcal F\circ\sigma$ حيث $\sigma f =
f(-\cdot)$، وهو
يحفظ الجداءات السلّمية إلى غاية العامل $2\pi$.
\end{theorem}

\begin{proof}
لتكن $f \in L^1\cap L^2$ ولتكن $h = f * \tilde f$ مع $\tilde f(x) = \overline{f(-x)}$. عندئذٍ $h \in L^1$
(\cref{thm:b3:product:convolution})، و $h$ متصلة ومحدودة (\cref{exo:b3:lp:6}: $f, \tilde f \in L^2$)،
و $h(0) =
\int f\bar f = \norm f_2^2$، و $\hat h = \hat
f\,\widehat{\tilde f} = \hat f\,\overline{\hat f} =
\abs{\hat f}^2 \geq 0$ (بحساب $\widehat{\tilde f} =
\overline{\hat f}$). ونطبّق \cref{thm:b3:fouriertransform:inversion}(1) على
$h$ عند $x = 0$:
\[
(h * g_\varepsilon)(0) = \frac1{2\pi}\int \hat
h(\xi)\,\eu^{-\varepsilon\xi^2}\dd\xi .
\]
وحين $\varepsilon \to 0$: يؤول الطرف الأيسر إلى $h(0)$ (لأن $h$ متصلة
محدودة: $(h*g_\varepsilon)(0) - h(0) = \int(h(-y)
- h(0))g_\varepsilon(y)\dd y \to 0$ بالشطر إلى $y$ صغيرة وكبيرة)؛ ويتزايد الطرف
الأيمن إلى $\frac1{2\pi}\int\hat h$ بمبرهنة التقارب الرتيب (لأن $\hat h \geq 0$).
ومنه $\frac1{2\pi}\int\abs{\hat f}^2 =
\norm f_2^2$، وهو منتهٍ أو غير منتهٍ \emph{قبليًّا} --- وهو
منتهٍ، فيَثبُت الانتماء والمتطابقة معًا.

وأما التمديد: فإن $L^1\cap L^2 \supseteq \mathcal C_c$ كثيف في $L^2$ (\cref{thm:b3:lp:density})؛ والتحويل
متقايس هناك إلى غاية العامل $\sqrt{2\pi}$، ومنه يمتدّ بكيفية
وحيدة إلى تقايس-إلى-غاية-ثابت $\mathcal F$ على $L^2$
(\cref{thm:b3:complete:extension}). وينتقل القلب من أجل $\mathcal S$ (\cref{thm:b3:fouriertransform:schwartz})
بالكثافة نفسها (إذ الطرفان متصلان بالمعنى $L^2$): أي $\mathcal
F\bigl(\frac1{2\pi}\mathcal F(\sigma f)\bigr) = f$
على $\mathcal S$، ومنه على $L^2$: فيكون تقابلًا. وأما الجداءات
السلّمية: فبالاستقطاب انطلاقًا من متطابقة المعيار.
\end{proof}

\begin{theorem}[صيغة بواسون في الجمع]
\label{thm:b3:fouriertransform:poisson}
لتكن $f \in \mathcal S(\R)$ (ويكفي أن تكون $f$ متصلة مع $\abs{f} +
\abs{\hat f} \leq C(1 + \abs\cdot)^{-2}$). عندئذٍ\index{صيغة
بواسون في الجمع}
\[
\sum_{n\in\Z} f(n) \;=\; \sum_{k\in\Z}\hat f(2\pi k) .
\]
\end{theorem}

\begin{proof}
لتكن $F(x) = \sum_{n\in\Z}f(x + n)$: تتقارب المتسلسلة تقاربًا ناظميًّا على المتراصات
(بتناقص $f$)، ومنه فإن $F$ متصلة، وهي دورية بالدور $1$. وأما
معاملات فورييه لها (بالدور $1$: $c_k(F) = \int_0^1F(t)\eu^{-2\iu\pi kt}\dd t$):
\[
c_k(F) = \sum_n\int_0^1 f(t + n)\,\eu^{-2\iu\pi kt}\dd t
= \int_\R f(t)\,\eu^{-2\iu\pi kt}\dd t = \hat f(2\pi k)
\]
(ويبرّر التقارب الناظمي التبديلَ؛ والطور دوري بالدور $1$).
وتتقارب المتسلسلة $\sum_k\abs{c_k(F)}$ (بتناقص $\hat f$)، ومنه تتقارب
متسلسلة فورييه للدالة $F$ تقاربًا ناظميًّا؛ ومجموعها دالة متصلة
لها معاملات فورييه نفسها التي للدالة $F$، ومنه تساوي $F$
(بالتباين على الدائرة: إذ للفرق معاملات معدومة، وتعطي
\cref{thm:b3:hilbert:fourier} أنه معدوم في $L^2$، ومنه في كل مكان بالاتصال). ثم
نقيّم عند $x = 0$.
\end{proof}

\begin{example}[متطابقة تيتا]
\label{ex:b3:fouriertransform:theta}
بتطبيق بواسون على $f(x) = \eu^{-\pi tx^2}$ (حيث $t > 0$)، وتحويلها $\hat f(\xi) = t^{-1/2}\eu^{-\xi^2/4\pi t}$
(\cref{ex:b3:fouriertransform:gaussian} مع $a = \pi t$):
\[
\sum_{n\in\Z}\eu^{-\pi n^2t}
= \frac1{\sqrt t}\sum_{k\in\Z}\eu^{-\pi k^2/t} :
\]
وهي المعادلة الدالية لدالة تيتا لياكوبي، ومفتاح المعادلة
الدالية للدالة $\zeta$ لريمان --- ومسرِّعٌ عددي باهر: فمن أجل
$t$ صغير، يتقارب الطرف الأيسر ببطء، ويتقارب الأيمن بسرعة
خاطفة.
\end{example}

\begin{method}\label{met:b3:fouriertransform:method}
مجالات العمل: $L^1$ --- حيث التحويل معرَّف نقطةً نقطة، والقلب
يحتاج إلى $\hat f \in L^1$؛ و $\mathcal S$ --- حيث كل شيء
مشروع، فابرهن هنا أولًا؛ و $L^2$ --- حيث التحويل معرَّف
بالكثافة (لا بالتكامل!)، بتناظر تام ومحاسبة بارسيفالية. ولحساب
تحويل: أرجعه إلى الجدول (الدالة المميّزة، والأُسّي، والغاوسية)
بقواعد \cref{prop:b3:fouriertransform:basic}؛ ولتبرهن على متطابقة: أثبتها على
$\mathcal S$ (أو $\mathcal C_c^\infty$) ومدّدها بالكثافة والاتصال
(\cref{met:b3:lp:toolkit})؛ ولتحل معادلة تفاضلية جزئية أو عادية خطية ذات
معاملات ثابتة: حوّل، واقسم، واقلب.
\end{method}

\begin{omfigure}
\begin{tikzpicture}[xscale=1.05, yscale=2.2]
  \draw[->] (-4.3,0) -- (4.3,0) node[right,font=\small] {$x$};
  \draw[->] (0,0) -- (0,1.25);
  \draw[omThm, thick, smooth]
    plot[domain=-4:4, samples=81]
    (\x, {exp(-\x*\x/0.4)/1.1209});
  \draw[omProp, thick, smooth]
    plot[domain=-4:4, samples=81]
    (\x, {exp(-\x*\x/1.6)/2.2418});
  \draw[omDef, thick, smooth]
    plot[domain=-4:4, samples=81]
    (\x, {exp(-\x*\x/6.4)/4.4836});
  \node[omThm, font=\small] at (1.35,0.75) {$t = 0.1$};
  \node[omProp, font=\small] at (2.4,0.36) {$t = 0.4$};
  \node[omDef, font=\small] at (3.5,0.17) {$t = 1.6$};
\end{tikzpicture}

{\small \omterm{pb:b3:fouriertransform:1}{نواة الحرارة} $g_t(x) =
\frac1{2\sqrt{\pi t}}\,\eu^{-x^2/4t}$ عند ثلاثة أزمنة: كتلة كلية $1$
إلى الأبد، وارتفاع $\sim t^{-1/2}$، وعرض $\sim
\sqrt t$. والتفاف
المعطيات الابتدائية مع هذه الغاوسية المنتشرة هو كل مضمون مسألة
نهاية الأسبوع؛ وفي مجال الترددات تُقرأ الصورة نفسها
$\hat g_t(\xi) = \eu^{-t\xi^2}$ --- فالترددات العالية تموت أولًا، وذلك اللاتناظر هو سهم
الزمن.}
\end{omfigure}

\section{تمارين}

\begin{exercise}[$\star$]\label{exo:b3:fouriertransform:1}
احسب تحويلات فورييه للدوال: $\mathbf 1_{\intcc{-a}a}$؛ و $\eu^{-a\abs x}$ (حيث
$a > 0$)؛ ودالة الخيمة $\max(0, 1 -
\abs x)$؛ و $\frac1{x^2 + a^2}$ \emph{(باستعمال
القلب على الثانية)}. وسجّل الجدول الناشئ.
\end{exercise}

\begin{exercise}[$\star$]\label{exo:b3:fouriertransform:2}
لتكن $f \in L^1$. عبّر بدلالة $\hat f$ عن تحويلات: $f(x - a)$،
و $f(x)\cos(bx)$، و $f(ax + b)$، و $\overline{f(-x)}$، و $(f * f)(x)$.
وتحقق من كل قاعدة على الدالة الغاوسية.
\end{exercise}

\begin{exercise}[$\star\star$]\label{exo:b3:fouriertransform:3}
(a) برهن على أن للمقدار $\mathbf 1_{\intcc{-1}1} * \mathbf
1_{\intcc{-1}1}$ التحويلَ $\bigl(\frac{2\sin\xi}\xi
\bigr)^2$، واستنتج
$\int_\R\bigl(\frac{\sin\xi}\xi\bigr)^2\dd\xi = \pi$ ببلانشرِل --- أو بالقلب عند $0$. وقارن \cref{pb:b3:lebesgue:1}.
(b) احسب $\int_\R\frac{\dd x}{(x^2+1)^2}$ ببلانشرِل مطبَّقة على $\eu^{-\abs x}$.
\end{exercise}

\begin{exercise}[$\star\star$]\label{exo:b3:fouriertransform:4}
(جبر \omterm{pb:b3:fouriertransform:1}{نواة الحرارة}) مع $g_t(x) =
\frac1{2\sqrt{\pi t}}\eu^{-x^2/4t}$:
(a) تحقق من $\hat g_t(\xi) = \eu^{-t\xi^2}$؛
(b) استنتج قانون نصف الزمرة $g_t * g_s = g_{t+s}$ دون أي حساب تكاملي؛
(c) برهن على $\norm{g_t}_1 = 1$ وعلى $\norm{g_t}_2^2 =
(8\pi t)^{-1/2}$.
\end{exercise}

\begin{exercise}[$\star\star$]\label{exo:b3:fouriertransform:5}
برهن على أنه إذا كانت $f \in L^1$ زوجية وحقيقية، كانت $\hat f$
زوجية وحقيقية؛ وإذا كانت $f$ فردية وحقيقية، كانت $\hat f$
فردية وتخيّلية محضة. وماذا يحسب $\hat f(0)$؟ واستنتج أن $f \geq
0$
يفرض $\norm{\hat f}_\infty = \hat f(0) = \int f$، وفسّر ذلك من أجل الكثافات الاحتمالية (\cref{ch:b3:clt}:
إذ لدالة مميِّزة طويلةٌ $\leq 1$، مبلوغة عند $0$).
\end{exercise}

\begin{exercise}[$\star\star\star$]\label{exo:b3:fouriertransform:6}
(عدم الغمر) برهن على أن $\widehat{\phantom f}\colon L^1
\to \mathcal C_0$ متباين ومتصل، لكنه \emph{غير}
غامر، في ثلاث خطوات.
(أ) التباين (\cref{thm:b3:fouriertransform:inversion}) والاتصال ($\norm{\hat f}_\infty \leq \norm f_1$)، وأن
$\mathcal C_0$ فضاء باناخ (لأنه مغلق في $\norm\cdot_\infty$).
(ب) لو كان التطبيق غامرًا لكان تقابلًا، ولأعطت مبرهنة التطبيق
المفتوح (\cref{thm:b3:banach:openmapping}) ثابتًا $C$ يحقق $\norm f_1 \leq C\norm{\hat
f}_\infty$ من أجل كل
$f \in L^1$.
(ج) ناقض ذلك بالمقدار $f_n(x) = \frac{\sin
x}{x}\cdot\frac{\sin(x/n)}{x/n}$: فتحويله (إلى غاية ثوابت) هو شبه
المنحرف من نمط الالتفاف $\mathbf 1_{\intcc{-1}1} *
\mathbf 1_{\intcc{-1/n}{1/n}}$ --- برهن على أن $\norm{\hat f_n}_\infty \leq \pi$
بانتظام، بينما $\norm{f_n}_1 \geq c\ln n$ بعدّ أقواس $\frac{\abs{\sin x}}x$ على $[1, n]$ (حيث
يكون العامل الثاني محدودًا من الأسفل)، كما في \cref{thm:b3:banach:fourierdiverge}.
\end{exercise}

\begin{exercise}[$\star\star$]\label{exo:b3:fouriertransform:7}
(معجم النعومة $\leftrightarrow$ التناقص) برهن على: أن $f \in
L^1$ مع $\hat f(\xi) = O(\abs\xi^{-k-1-\delta})$
من أجل $\delta > 0$ ما يستلزم أن تقبل $f$ ممثّلًا من الصنف
$\mathcal C^k$. وبالعكس، يستلزم $f \in \mathcal C^k_c$ أن $\hat f(\xi) =
O(\abs\xi^{-k})$. ووضّح
الاتجاهين على دالة الخيمة.
\end{exercise}

\begin{exercise}[$\star\star\star$]\label{exo:b3:fouriertransform:8}
(متراجحة هايزنبرغ) من أجل $f \in \mathcal S(\R)$ حقيقية مع $\norm f_2 = 1$،
برهن على
\[
\Bigl(\int x^2f(x)^2\dd x\Bigr)\cdot
\Bigl(\frac1{2\pi}\int \xi^2\abs{\hat f(\xi)}^2\dd\xi\Bigr)
\;\geq\; \frac14 ,
\]
مع التساوي من أجل الدوال الغاوسية. \emph{(اكتب $1 = \int f^2 =
-\int x\,(f^2)'$
بالتجزئة، وحُدّه بكوشي--شوارتز، وحوّل $\norm{f'}_2$
ببلانشرِل.)} والتفسير: أن إشارةً وطيفَها لا يمكن أن يتركّزا
معًا.
\end{exercise}

\begin{exercise}[$\star\star$]\label{exo:b3:fouriertransform:9}
برّر \cref{ex:b3:fouriertransform:theta} بالتفصيل (أي فرضيات بواسون من أجل الدالة
الغاوسية)، واستعمل المتطابقة لتقييم $\sum_{n\in\Z}\eu^{-\pi n^2}$ بستة أرقام عشرية
بثلاثة حدود. وكم حدًّا من المتسلسلة \emph{المعرِّفة} تلزم لبلوغ
الدقة نفسها عند $t = 10^{-2}$، مقابل المتسلسلة المحوَّلة؟
\end{exercise}

\begin{exercise}[$\star\star$]\label{exo:b3:fouriertransform:10}
(الدوال المحدودة الحزمة) لتكن $f \in L^2(\R)$ بحيث يكون حامل
$\mathcal Ff$ في $\intcc{-\pi}\pi$. برهن على أن $f$ تقبل ممثّلًا يمتدّ
بحيث تكون كل قيمة من قيمه قابلة للاستعادة من العيّنات: أي برهن
على \emph{استيفاء شانون} عند الأعداد الصحيحة،
\[
f(x) = \sum_{n\in\Z} f(n)\,
\frac{\sin\bigl(\pi(x - n)\bigr)}{\pi(x - n)}
\quad\text{في} L^2,
\]
بنشر $\mathcal Ff$ في أساس فورييه للفضاء $L^2(\intcc{-\pi}\pi)$ (\cref{thm:b3:hilbert:fourier})
ثم بالتحويل العكسي حدًّا حدًّا.
\end{exercise}

\begin{exercise}[$\star\star$]\label{exo:b3:fouriertransform:11}
(التحويل بوصفه مؤثرًا من الرتبة أربعة) على $\mathcal
S(\R)$، لتكن
$\mathcal F f = \hat f$.
(a) باستعمال صيغة القلب، برهن على $(\mathcal F^2f)(x) =
2\pi\,f(-x)$، واستنتج $\mathcal F^4 = (2\pi)^2\,
\mathrm{id}$.
(b) استنتج أن كل قيمة ذاتية للمؤثر $\mathcal F$ على
$\mathcal S$ تنتمي إلى $\{\pm\sqrt{2\pi},
\pm\iu\sqrt{2\pi}\}$، وأبرز دالة ذاتية موافقة
للقيمة $+\sqrt{2\pi}$ \emph{(أي دالة من دوال هذا الفصل تتناسب
مع تحويلها؟)}.
(c) برهن على أن الدوال الزوجية تحقق $\mathcal F^2f = 2\pi
f$ والفردية تحقق
$\mathcal F^2f = -2\pi f$؛ وأنتج دالة ذاتية للقيمة الذاتية $-\iu\sqrt{2\pi}$ انطلاقًا من
$x\eu^{-x^2/2}$ بحساب تحويلها (باشتقاق تحويل الدالة
الغاوسية).
\end{exercise}

\begin{exercise}[$\star\star$]\label{exo:b3:fouriertransform:12}
(الارتباط الذاتي والمبرهنة المساعدة لفينر) من أجل
$f \in L^2(\R)$ نعرّف $\tilde f(x) = \overline{f(-x)}$ و\emph{الارتباط الذاتي} $A_f = f * \tilde f$.
(a) برهن على أن $A_f$ دالة متصلة محدودة تحقق $A_f(0) = \norm f_2^2 \geq \abs{A_f(x)}$ من أجل كل
$x$ (\cref{exo:b3:lp:6} وكوشي--شوارتز).
(b) برهن، أولًا من أجل $f \in L^1\cap L^2$، على أن $\widehat{A_f} = \abs{\hat f\,}^2 \geq 0$: أي إن للارتباط
الذاتي تحويلًا غير سالب --- فأطياف الارتباطات الذاتية أطيافُ
استطاعة.
(c) استنتج المتطابقة $\int_\R\abs{\hat
f(\xi)}^2\eu^{\iu x\xi}\,\dd\xi = 2\pi A_f(x)$ (بالقلب؛ وبرّر انطباقه حين يحقق
$\hat f \in L^2$ الشرطَ $\abs{\hat f}^2 \in L^1$)، وقيّمها من أجل $f = \mathbf
1_{\intcc{-1/2}{1/2}}$ عند
$x = 0$: لتستعيد $\int_\R\bigl(\frac{\sin u}u\bigr)^2\dd u = \pi$.
\end{exercise}

\section{مسألة: معادلة الحرارة على المستقيم}

\begin{problem}[{مسألة نهاية الأسبوع --- $\partial_tu =
\partial^2_{xx}u$، محلولةً من
أولها إلى آخرها}]
\label{pb:b3:fouriertransform:1}
تنتشر الحرارة؛ وتقول المعادلة $\partial_tu = \partial_{xx}^2u$ إن كثافتها تنتثر بمعدل
يعطيه التقوّس المحلي لملمح درجة الحرارة. ونحل مسألة كوشي على
$\R$ --- أي بإعطاء $f$، نجد $u(t, x)$ من أجل $t > 0$ يحقق
$u(0, \cdot) = f$ --- ونبرهن على خصائص الحل اللافتة، ونرى لماذا لا يمكن
عكس الزمن. وفي كل ما يلي، تكون $g_t(x) = \frac1{2\sqrt{\pi t}}\eu^{-x^2/4t}$ \emph{نواة
الحرارة}\index{نواة الحرارة} وتكون $u(t, \cdot) = g_t *
f$.

\medskip
\noindent\textbf{الجزء الأول --- استنباط النواة.} لنعمل شكليًّا
أولًا: لنفترض أن $u(t, \cdot) \in \mathcal S$ تحل المعادلة، وليكن $\hat u(t, \xi)$
التحويلَ في $x$.
\begin{enumerate}
  \item برهن (شكليًّا) على $\partial_t\hat u = -\xi^2\hat u$، ومنه $\hat u(t,\xi) = \eu^{-t\xi^2}\hat f(\xi)$، وتعرّف على
        $u(t) = g_t * f$ (\cref{exo:b3:fouriertransform:4}). وهذا يحفّز
        \emph{تعريف} $u$؛ وكل شيء يُبرهَن عليه الآن مباشرةً،
        من أجل $f \in \mathcal C_b(\R)$ (المحدودة المتصلة) أو $f \in L^p$.
\end{enumerate}

\medskip
\noindent\textbf{الجزء الثاني --- التحقق.}
\begin{enumerate}[resume]
  \item برهن على أنه من أجل $t > 0$، تكون $u(t, x) = \int
        g_t(x-y)f(y)\dd y$ معرَّفة
        جيدًا من أجل $f \in \mathcal
        C_b$، وعلى أن $u$ من الصنف $\mathcal C^\infty$ في
        $(t, x)$ على $\intoo0\infty\times\R$ \emph{(بالاشتقاق تحت علامة
        التكامل؛ وبالهيمنة على المشتقات الغاوسية بانتظام
        محليًّا في $(t,x)$)}.
  \item تحقق من $\partial_tg_t = \partial^2_{xx}g_t$ بالحساب المباشر، واستنتج $\partial_tu =
        \partial^2_{xx}u$ من
        أجل $t > 0$.
  \item (الشرط الابتدائي) برهن على أنه من أجل $f \in \mathcal
        C_b$، يكون
        $u(t, x) \to f(x)$ حين $t \to 0^+$، بانتظام على المتراصات
        \emph{(بالوحدة التقريبية: اشطر $\abs y
        \leq \delta$، $\abs y > \delta$)}؛ ومن
        أجل $f \in L^p$ (حيث $p < \infty$)، برهن على $\norm{u(t) - f}_p \to 0$.
  \item (التنعيم الفوري) اخلص: أنه حتى من أجل $f$ محدودة
        متصلة فحسب، يكون الحل $\mathcal C^\infty$ من أجل كل $t > 0$ ---
        فالحرارة تمحو الخشونة فورًا. واحسب $u(t, \cdot)$
        صراحةً من أجل $f =
        \mathbf 1_{\intoo0\infty}$ (وهي دالة خطأ)، وارسم ملمحه من
        أجل ثلاث قيم للزمن $t$.
\end{enumerate}

\medskip
\noindent\textbf{الجزء الثالث --- الخصائص البنيوية.}
\begin{enumerate}[resume]
  \item (الإيجابية والمقارنة) إذا كانت $f \geq 0$ فإن $u > 0$
        من أجل كل $t > 0$ (تمامًا، إلا إذا كانت $f = 0$ في كل
        مكان تقريبًا)؛ وإذا كانت $f_1 \leq f_2$ فإن
        $u_1 \leq u_2$. فبقعةٌ باردة تدفأ فورًا: علّق على
        ذلك.
  \item (الحفظ) من أجل $f \in L^1$: $\int u(t, x)\dd x =
        \int f$ من أجل كل $t$
        \emph{(بتونيلي)} --- فالحرارة الكلية محفوظة.
  \item (التبديد) من أجل $f \in L^1\cap L^2$، برهن ببلانشرِل على أن
        $t \mapsto \norm{u(t)}_2$ متناقصة بالمعنى الواسع، وتمامًا إلا إذا كانت
        $f = 0$، واحسب نهايتها حين $t \to \infty$. وبرهن فوق
        ذلك على $\norm{u(t)}_\infty \leq \frac{\norm
        f_1}{2\sqrt{\pi t}} \to 0$: فالحرارة تنتشر وتموت.
  \item (الوحدانية، الصنف $L^2$) ليكن $u$ حلًّا يحقق
        $u(t) \in L^2$ من أجل كل $t$، و $u \in \mathcal
        C^1(\intoo0\infty, L^2)$ بالمعنى
        الطبيعي، و $u(t) \to f$ في $L^2$ حين $t\to0$؛
        ومسلِّمًا بأن التحويل يحوّلها إلى $\partial_t\hat u =
        -\xi^2\hat u$ نقطةً نقطة
        في كل مكان تقريبًا في $\xi$ من أجل $t$ في كل مكان
        تقريبًا \emph{(وهو مبرَّر بالاختبار مقابل $\mathcal
        C_c^\infty$ في
        $\xi$ --- اعرض ذلك في خطوطه العامة)}، برهن على
        $\hat u(t,\xi) = \eu^{-t\xi^2}\hat f(\xi)$، ومن ثَمّ على الوحدانية في هذا الصنف.
\end{enumerate}

\medskip
\noindent\textbf{الجزء الرابع --- سهم الزمن.}
\begin{enumerate}[resume]
  \item برهن على أن المسألة \emph{العكسية} سيّئة الوضع: فلكي
        يوجد الحل عند الزمن $-s$ (حيث $s > 0$) بالمعطيات $f$
        عند الزمن $0$ --- أي لكي تقبل $f = g_s * h$ حلًّا
        $h \in L^2$ --- يلزم أن يكون $\eu^{s\xi^2}\hat f(\xi) \in L^2$: وهو شرط تناقص
        متطرّف على $\hat f$. وأبرز دالةً ناعمة صريحة $f
        \in L^2$
        لا يوجد من أجلها أي حل عكسي على أي فترة زمنية: خذ
        الدالة التي $\hat f(\xi) =
        \eu^{-\abs\xi}$ --- وعيّن $f$ (\cref{exo:b3:fouriertransform:1}) وبرهن
        على $\eu^{s\xi^2}\eu^{-\abs\xi} \notin L^2$ من أجل كل $s >
        0$.
  \item (التنعيم مقابل المعلومة) اشرح في فقرة قصيرة،
        باستعمال الأسئلة 5 و9 و10، لماذا تكون نصف زمرة
        الحرارة $(f \mapsto g_t * f)_{t\geq0}$ متباينةً وغير غامرة على $L^2$، ولماذا
        يعبّر ذلك عن لاعكوسية الانتثار.
\end{enumerate}

\medskip
\noindent\textbf{الجزء الخامس --- مبرهنة شانون في أخذ
العيّنات.} تكون الدالة $f \in L^2(\R)$ \emph{محدودة الحزمة}
بالمقدار $\Omega$ إذا كان $\hat f = 0$ في كل مكان تقريبًا خارج
$\intcc{-\Omega}\Omega$؛ ونكتب $PW_\Omega$ (فضاء بالي--فينر) لهذه الدوال.
\begin{enumerate}[resume]
  \item برهن على أن كل $f \in PW_\Omega$ توافق في كل مكان تقريبًا
        الدالةَ $\mathcal C^\infty$ من الصنف $\frac1{2\pi}\int_{-\Omega}^{\Omega}\hat
        f(\xi)\eu^{\iu x\xi}\,\dd\xi$ (وبرّر النعومة
        والتعيين)، وأن جميع مشتقاتها محدودة: فتحديد الحزمة
        صورةٌ متطرّفة من الانتظام. ومن الآن فصاعدًا ترمز $f$
        إلى هذا الممثّل.
  \item انشر $\hat f \in L^2(\intcc{-\Omega}\Omega)$ في أساس فورييه لتلك الفترة وعيّن
        المعاملات بوصفها \emph{عيّنات} للدالة $f$:
        \[
        \hat f(\xi) = \frac\pi\Omega\sum_{n\in\Z}
        f\Bigl(\frac{n\pi}\Omega\Bigr)\,
        \eu^{-\iu n\pi\xi/\Omega}
        \quad\text{في} L^2(\intcc{-\Omega}\Omega) .
        \]
  \item استنتج \emph{مبرهنة أخذ العيّنات}: أنه من أجل $f \in
        PW_\Omega$،
        \[
        f(x) = \sum_{n\in\Z}f\Bigl(\frac{n\pi}
        \Omega\Bigr)\,\operatorname{sinc}(\Omega x - n\pi),
        \qquad \operatorname{sinc}t = \frac{\sin t}t,
        \]
        مع التقارب في $L^2(\R)$ وبانتظام على $\R$ \emph{(أدخل
        متسلسلة السؤال 13 في صيغة القلب واحسب التكامل
        الابتدائي)}: فالإشارة المحدودة الحزمة معيَّنة تمامًا
        بقيمها على شبكة خطوتها $\pi/\Omega$ --- أي معدل
        نايكويست.
  \item برهن على أن الدوال $x \mapsto
        \operatorname{sinc}(\Omega x - n\pi)$، حيث $n \in \Z$، تشكّل
        عائلة متعامدة في $L^2(\R)$ بمعيار ثابت $\sqrt{\pi/\Omega}$،
        واستنتج متطابقة الطاقة $\norm f_2^2 =
        \frac\pi\Omega\sum_n\abs{f(n\pi/\Omega)}^2$.
  \item (الالتباس الطيفي) أبرز دالة غير معدومة $g \in PW_{2\Omega}$ تنعدم
        عند كل نقطة عيّنة $\frac{n\pi}\Omega$ \emph{(انظر في $g(x) =
        \sin(\Omega x)\operatorname{sinc}(\Omega x)$
        وتحقق من حزمتها)}: فأخذ العيّنات دون معدل نايكويست
        يفقد المعلومة --- إذ يمكن لإشارتين مختلفتين أن
        تتقاسما جميع العيّنات: أي أثر عجلة العربة
        الستروبوسكوبي، مصوغًا رياضيًّا.
  \item (درجات الحرية) باستعمال السؤالين 14 و15، برّر القاعدة
        الهندسية: أن إشارةً محدودة الحزمة بالمقدار $\Omega$
        تحمل نافذةٌ زمنية طولها $T$ طاقتَها جوهريًّا، تُوصَف
        بنحو $\frac{\Omega T}\pi$ عيّنة حقيقية --- واضبط معنى «جوهريًّا»
        بواسطة متطابقة الطاقة والذيل $\sum_{\abs{n\pi/\Omega} >
        T/2}$.
  \item (فحوص الاتساق) تحقق من مبرهنة أخذ العيّنات يدويًّا
        على عنصرين من $PW_\Omega$: (a) $f =
        \operatorname{sinc}(\Omega\,\cdot)$، وعيّناتها
        $\delta_{n0}$؛ (b) الإشارات الضيقة الحزمة من نمط
        $f(x) = \cos(\omega x)
        \operatorname{sinc}(\varepsilon x)$ --- وبمزيد من الدقة، برهن على أنه من أجل
        $f \in
        PW_{\Omega'}$ مع $\Omega' < \Omega$، تستعيد متسلسلة المعدل $\Omega$
        الدالةَ $f$ أيضًا (فالإفراط في أخذ العيّنات غير ضار)،
        بغمر $PW_
        {\Omega'} \subseteq PW_\Omega$.
\end{enumerate}

\medskip
\noindent\textbf{الجزء السادس --- اللايقين، مرتين أخريين.}
تحدّ متراجحة هايزنبرغ (\cref{exo:b3:fouriertransform:8}) \emph{مقدار} ما يمكن أن
تتركّز به $f$ و $\hat f$ معًا؛ وهذان شقيقتها من نمط الكل أو
لا شيء وإشباعها المضبوط.
\begin{enumerate}[resume]
  \item لتكن $f \in L^1$ مع $\operatorname{supp}f \subseteq
        \intcc{-A}A$. برهن على أن $\hat f$
        مجموع متسلسلة قوى متقاربة في كل مكان:
        \[
        \hat f(\xi) = \sum_{k\geq0}\frac{(-\iu\xi)^k}{k!}
        \,m_k, \qquad m_k = \int_{-A}^{A}x^kf(x)\,\dd x,
        \quad \abs{m_k} \leq A^k\norm f_1
        \]
        \emph{(انشر $\eu^{-\iu\xi x}$ وبرّر التبديل بالتقارب
        الناظمي)}: فالتحويل \emph{تحليلي حقيقي}، بنصف قطر
        تقارب لانهائي عند كل نقطة.
  \item استنتج ثنائية الحامل: أن دالة تحليلية حقيقية تنعدم
        على فترة مفتوحة غير خالية تنعدم انعدامًا تامًّا
        \emph{(فالمجموعة التي تنعدم عندها جميع المشتقات
        مفتوحة ومغلقة --- فصّل حجة تايلور)}؛ واخلص إلى أنه لا
        توجد $f$ غير معدومة يكون حاملُ كلٍّ من $f$ و
        $\hat f$ متراصًّا، وإلى أن $PW_\Omega$ لا يحتوي أي
        دالة غير معدومة ذات حامل متراص --- فالإشارات المحدودة
        الحزمة تدوم إلى الأبد، والإشارات المحدودة الزمن
        تتسرّب إلى جميع الترددات.
  \item (إشباع هايزنبرغ) على العائلة الغاوسية $f
        = \eu^{-ax^2}$، احسب
        عاملَي التركّز كليهما وتحقق من أن الجداء المعيَّر
        $\bigl(\int x^2\abs f^2\bigr)\bigl(\frac1{2\pi}\int
        \xi^2\abs{\hat f}^2\bigr)\big/\norm f_2^4$ يساوي $\frac14$ من أجل \emph{كل} $a$ --- أي
        عائلة التساوي في \cref{exo:b3:fouriertransform:8} بلحمها ودمها؛ واشرح بحجة
        تحجيمية لماذا يجب أن يكون الجداء ثابتًا على امتداد
        العائلة.
  \item اشرح القراءة الفيزيائية (كثافتا الموضع والاندفاع
        لحالة كمومية؛ ويعطي $\hbar$ في التعيير $\sigma_x\sigma_p \geq
        \frac\hbar2$)، وصِل
        بين أجزاء الفصل: فالتنعيم الفوري (الجزء الثاني)،
        واللاعكوسية (الجزء الرابع)، وأخذ العيّنات (الجزء
        الخامس)، وهايزنبرغ، وثنائية الحامل، خمسةُ تعبيرات عن
        قانون واحد --- أي إن سلوك $\hat f$ عند اللانهاية يشرّع
        ما يمكن أن تفعله $f$ في أي مكان.
\end{enumerate}

\medskip
\noindent\textbf{الجزء السابع --- جبر النواة، ومثال واحد قابل
للحل.}
\begin{enumerate}[resume]
  \item (نصف الزمرة) برهن على \emph{متطابقة
        تشابمان--كولموغوروف} $g_t * g_s = g_{t+s}$ من أجل $t, s > 0$ (عبر
        مبرهنة الالتفاف وتباين التحويل على $L^1$)، واستنتج
        $u(t + s) = g_s *
        u(t)$: أي إن التطوّر خلال الزمن $t + s$ هو التطوّر
        خلال $t$ ثم خلال $s$. وحدّد تبديد السؤال 8: فبكتابة
        $\norm{u(t)}_2^2 = \frac1{2\pi}\int
        \eu^{-2t\xi^2}\abs{\hat f(\xi)}^2\dd\xi$، برهن بكوشي--شوارتز على أن
        \[
        t \longmapsto \ln\,\norm{u(t)}_2
        \quad\text{محدّبة على} \intoo0{+\infty} :
        \]
        أي إن الطاقة $L^2$ لا تتناقص فحسب، بل تتناقص بكيفية
        محدّبة لوغاريتميًّا.
  \item (إلى أين تذهب الحرارة) لتكن $f \geq 0$ و $f \in L^1$،
        مع $\int x^2f(x)\dd x < \infty$. برهن على أنه من أجل كل $t > 0$
        \[
        \int_\R x\,u(t, x)\,\dd x = \int_\R x f(x)\,\dd x,
        \qquad
        \int_\R x^2u(t, x)\,\dd x
        = \int_\R x^2f(x)\,\dd x + 2t\int_\R f :
        \]
        أي إن مركز الحرارة لا يتحرك أبدًا، وأن التباين ينمو
        \emph{خطيًّا} في الزمن --- وهو التحجيم الانتثاري
        $x \sim \sqrt{2t}$، الذي يُعاد قراءته حين تظهر الحركة البراونية
        في \cref{ch:b3:probability}. \emph{(احسب العزمين الأولين للنواة
        $g_t$ واستعمل تونيلي على الالتفاف.)}
  \item (الدالة الغاوسية، محلولةً من أولها إلى آخرها) من أجل
        $f(x) =
        \eu^{-x^2}$، أثبت الصيغة المغلقة
        \[
        u(t, x) = \frac1{\sqrt{1 + 4t}}\,
        \exp\Bigl(-\frac{x^2}{1 + 4t}\Bigr),
        \]
        وتحقق عليها يدويًّا من: المعادلة $\partial_tu
        = \partial^2_{xx}u$؛ والحفظ
        $\int u(t) =
        \sqrt\pi$؛ وقانون التبديد $\norm{u(t)}_2 =
        (\pi/2)^{1/4}(1 + 4t)^{-1/4}$ (وقارن تناقصه
        $t^{-1/4}$ بتناقص معيار السوپريموم $t^{-1/2}$ في
        السؤال 8)؛ ونمو التباين المضبوط في السؤال 24. وعند
        $t = 6$: تكون الذروة قد هبطت إلى $\frac15$ من
        ارتفاعها الابتدائي بينما صار الملمح أعرض بخمسة أضعاف
        --- الحرارة نفسها، منتشرةً.
\end{enumerate}
\end{problem}
